(1) 光线从空气以轴向夹角 \(\theta_a\) 入芯，Snell定律 \(n_0\sin\theta_a=n_1\sin\theta_1\)。在芯包界面，全反射要求内部入射相对法线角不小于临界角，等价 \(\sin\theta_1\le\sqrt{1-(n_2/n_1)^2}\)。所以 \(\boxed{n_0\sin\theta_{a,\max}=\mathrm{NA}=\sqrt{n_1^2-n_2^2}}\)，实际空气入射角还受 \(\sin\theta\le1\) 限制。

(2) 芯内角满足 \(\theta_1\le\arccos(n_2/n_1)\)，弱导情况下 \(\theta_{1,\max}\simeq\sqrt{2\Delta}\) 弧度。这个结论来自几何全反射；精确电磁本征模还要满足圆柱界面的切向场连续与有限包层衰减，不能把每条射线都称作独立传输模。

(3) 定义 \(\boxed{V=(2\pi a/\lambda)\sqrt{n_1^2-n_2^2}}\)。标准弱导阶跃圆柱光纤中当 \(V<2.405\) 时只支持基模LP\(_{01}\) 的一组偏振态；其理想导模无有限正 \(V\) 截止，而更高横模的第一截止发生在 \(V\simeq2.405\)。

(4) \(V\propto a/\lambda\)，所以增大波长或减小芯半径使 \(V\) 下降、可支持的高阶导模逐渐减少，易进入单模区。几何光线条件只判定芯包反射可能性，没有圆柱横向本征量子化，所以单模阈值必须用波动本征问题而不是仅用NA计算。
